\pg{355} \emph{principio}, de que se parte e que importa alguma coisa, a qual se perde; --- 2º) um \emph{fermo}, a que se chega e que importa alguma coisa, a qual se. adquire; --- 3*) um \emph{sujeito}, que do \emph{principio} chega ao \emph{fermo} e se encontra num e noutro extremo. --- Se faltar algum d' estes elementos, não haverá verdadeira mudança \footnote{O \emph{principio}, de que se parte, é chamado, pelos escholasticos, termo '\& \emph{quo}; \emph{o termo}, a que se chega, termo \emph{dd} quem. --- Entre o \emph{principio} e o termo, ou entre o termo a \emph{quo} é o termo \emph{ad} quem, deve existir uma certa opposiqio; aliás não poderia conceber-se a mundauca. Diz S. Thomaz: « Omnis mutatio est inter opposita» (In V \emph{Physic}, 1. 9. --- A mudança recebe a especie e 2 denominação --- não do \emph{principio}, ou do termo a quo, --- mas do fermo, ou do termo ad quem, porque é para este termo que a mudança tende e é nelle que se completa; assim dizemos que a agua está quente, quando de fria se muda para quente. £screve o S. Doutor : « Mutaliones accipiunt speciem et dignitatem non a termino a \emph{quo}, sed a. termino \emph{ad} quem» (dum. \emph{Th}., p. I, q. 45, a. 1 ad 2). --- O \emph{sujeito} da mudança deve Ser um e o mesmo no \emph{principio} e no \emph{fermo}; o que muda propriamente é \emph{o estado} do sujeito. Esta mudança do estado realisa-se emquanto o sujeito perde Uma coisa e adquire outra, --- Se, pois, o sujeito não se encontrasse DO \emph{principio}, ou no \emph{fermo}, não haveria. propriamente \emph{mudança}. Assim não é mudança --- nem a creação, porque antes d' ella não existe o sujeito, isto é, não ba um sujeito commum aos dois extremos, --- nem a \emph{anniguilaczo}, porque depois d' ella não existe o sujeito, isto 6, não ha um sujeito commum aos dois extremos, pois o \emph{nada} não é um \emph{estado}, mas uma mera negação (C£ \emph{C}. \emph{Gent}., II, 17; \emph{De Pot}, q. 3, a. 2)}.

S£. Especies de mudança. --- A \emph{mudança} é --- \emph{intrinsect} ou \emph{extrinseca}, --- \emph{instantanea} ou successiva, --- \emph{substancial} ou \emph{accidental},

a) \emph{Mudança intrinseca} ou \emph{extrinseca}, --- É \emph{intrinseca}, quando o sujeito adquire ou perde alguma coisa em si mesmo; tal é a mudança de um corpo, que de frio se torna quente. --- É \emph{extrinseca}, quando o sujeito nada adquire ou perde em si mesmo, mas apenas reveste uma nova relação por se ter mudado um outro ente, que de algum modo lhe se refere; assim dizemos que uma torre, que estava dá nossa direita, se---rmada-em---- quanto passa para a nossa esquerda pelo facto. de termos. nós mudado de posição \footnote{Esta distineção não deve confundir-se com a distinecio da mudança proveniente de uma.causa \emph{inírinseca,O0u} de uma causa \emph{extrinseca}. --- Mudança proveniente de uma causa \emph{intrinseca} É a que se funda na propria natureza de sujeito, que tem a potencialidade para a não-existencia, como \&. toda a substancia composta de materia e de forma. --- Mudança proveniente de uma cansa \emph{extrinseca} é a que tem por \emph{principio} a potencia activa de \emph{um} agente exterior, de que o sujeito depende ; assim a nossa alma, embora não tenha\pgn{356}em si a potencia para a não-existencia, comtudo, por estar sujeita 2o poder omnipotente de Deus, pode mudar. --- Como se vé, a mudança \emph{intrinseca} pode derivar de uma causa \emph{exirinseca}.}

